% =====================================================================
% Response to the supervisor's second round of \federico{} comments on
% Paper D (v10 -> v11). Standalone: pdflatex twice.
% =====================================================================
\documentclass[11pt,a4paper]{article}
\usepackage[T1]{fontenc}
\usepackage[utf8]{inputenc}
\usepackage[expansion=false]{microtype}
\usepackage[margin=2.4cm]{geometry}
\usepackage{amsmath,amssymb}
\usepackage{enumitem}
\usepackage{xcolor}
\usepackage[hidelinks]{hyperref}
\definecolor{qgrey}{HTML}{4A4A4A}
\definecolor{rule}{HTML}{C9C9C4}
\definecolor{tag}{HTML}{8A4B2A}
\setlength{\parskip}{0.35em}
\setlength{\parindent}{0pt}
\newcommand{\thinrule}{\textcolor{rule}{\rule{\linewidth}{0.4pt}}}
\newenvironment{qitem}[3]{%
  \par\medskip\thinrule\par\smallskip
  \noindent\textbf{\textcolor{tag}{#1}}\quad\textcolor{qgrey}{\small #2}\par\smallskip
  \noindent\textcolor{qgrey}{\small\itshape #3}\par\smallskip
}{\par}
\newcommand{\where}[1]{\par\smallskip{\small\textcolor{tag}{$\hookrightarrow$}\ \textit{#1}}\par}

\begin{document}
\begin{center}
{\Large\bfseries Response to Supervisor Comments, Second Round}\\[0.35em]
{\large Paper D: \emph{From Wire to Flight: Certifying UAV Mission Completion
Against Detector-Evading Attacks}}\\[0.6em]
\end{center}

Thank you for the second reading. Each comment is answered below with the
place where the revision was made. Two of them changed more than wording. The
comment on the six datasets exposed an overstatement: one of the six was never
ingested. The suggested related work bears directly on our novelty claim, which
is now positioned against it. The paper remains within the S\&P limit
(13~pages of body, 18 in total).

\begin{qitem}{1}{Abstract, terminology and claims}{Use ``unmanned aerial
vehicle (UAV)'' once, then ``UAV''; reserve ``certified''; avoid ``none'',
``never'', ``first'', ``universal'', ``any'', ``no prior''; one term each for
corpus and benchmark; define acronyms once; the novelty claim is too strong.}
All six points are applied across the paper.
\begin{itemize}[leftmargin=1.4em,itemsep=1pt]
\item ``Unmanned aerial vehicle (UAV)'' is defined in the abstract and at its
first use in the body, and ``UAV'' is used thereafter. ``Drone'', ``aircraft''
and ``uncrewed'' no longer denote the vehicle; ``airframe'' remains only where it
means the physical platform.
\item The Scope paragraph now states: ``Throughout, `certified' means proved
under (C1)--(C3).'' The paper has three conditions, so we write (C1)--(C3)
rather than (C1)--(C4); please tell us if you had a fourth condition in mind.
\item Empirical and rhetorical absolutes are qualified (``in the settings we
simulate'', ``in our evaluation'', ``to the best of our knowledge'').
Mathematical quantifiers inside definitions, theorems and proofs are unchanged.
\item ``Corpus'' now means the six sources and ``benchmark'' the corpus, schema
and harness together.
\item Every acronym is defined at its first use in the body and not again. The
background appendix no longer redefines terms.
\item The novelty claim is positioned against the related work in item~3.
\end{itemize}
\where{Abstract, Introduction, Scope, throughout.}
\end{qitem}

\begin{qitem}{2}{Definition 1}{``Every attack'' is too broad, since the
adversary can also drop or forge traffic and manipulate RF.}
Agreed. The definition now reads ``for every time-delay attack $a$ with
$D_\theta(a)=$ clean'', which is the class the theorem covers.
\where{\S II-D, Definition 1.}
\end{qitem}

\begin{qitem}{3}{Related work}{Consider Mo \& Sinopoli (TAC 2016), Urbina et
al.\ (CCS 2016), Choi et al.\ (CCS 2018) and SAVIOR (USENIX Security 2020).}
Added, as a new paragraph, ``Impact of Stealthy Attacks''. These works share our
premise that a detector's threshold limits what an attacker can achieve, and
they establish it for integrity attacks within a single control loop, where the
detector observes the same physical residuals that the attack perturbs. We now
describe our contribution as transferring that premise across layers. Our
detector observes network forwarding, the attack delays messages rather than
falsifying them, the bounded quantity is the fraction of missions that complete,
and the link between the layers is measured rather than modelled. The
Introduction was adjusted so that it does not suggest that no one has linked a
threshold to bounded damage.
\where{\S III (new paragraph), Introduction.}
\end{qitem}

\begin{qitem}{4}{Appendix references}{Every reference to an appendix section
prints ``Appendix A'', and the appendix itself is just called ``Appendix''.}
This was a LaTeX error (\verb|\appendix| instead of IEEEtran's
\verb|\appendices|). The appendices are now lettered A--G, and every reference
resolves to its own letter.
\where{Appendices.}
\end{qitem}

\begin{qitem}{5}{Corpus table}{You talk about six datasets, but are you using
all of them in the experiments?}
No, and the paper now says so. Adapters exist for all six, but only five are
ingested: the UAV-CAS release could not be retrieved for this study, so only its
adapter is included. The ingestion table had already listed it as ``adapter
ready; file pending'', while the abstract and introduction said ``we unify six''.
They now say ``adapters for six, five ingested''. The corpus section and the
table caption also state each dataset's role:
\begin{itemize}[leftmargin=1.4em,itemsep=1pt]
\item Three supply the certificate's constants: DATAMUt (residual budget and
operating curve), the UAV Attack Dataset (the rate $\gamma$) and UAV-EW-Bench
(the completion anchor).
\item UAVIDS-2025 and HCRL UAVCAN exercise the schema and contribute the corpus
statistics and the bus-attack classes, but they do not enter the certified
chain.
\end{itemize}
Paper IV made the same claim about UAV-CAS and is corrected in the same way.
\where{Abstract, Introduction, \S V and the Table I caption; related work no
longer says ``we ingest'' UAV-CAS.}
\end{qitem}

\begin{qitem}{6}{Proposition 1 (Soundness)}{There are no conditions in the
statement, but the proof relies on (C1)--(C3).}
Agreed. The proposition now begins ``Under conditions (C1)--(C3)'' and covers
every $D_\theta$-evading time-delay attack.
\where{\S VIII, Proposition 1.}
\end{qitem}

\begin{qitem}{7}{Figure (residual budget)}{``Scenario~1'' is never defined.}
The harness section now defines the three scenarios: scenario~1 is a $4\times4$
grid, scenario~2 a nine-node ring and scenario~3 a double ring. The caption
names the grid.
\where{\S VI and the residual-budget figure caption.}
\end{qitem}

\begin{qitem}{8}{Appendix on learned detectors}{$\sigma$ is already used for the
smoothing noise.}
The learned detector's score is now $\zeta(a)$, so $\sigma$ denotes only the
smoothing noise.
\where{Appendix D.}
\end{qitem}

\end{document}
